\documentclass[12pt,a4paper]{article}
\usepackage{amsmath, amssymb, amsthm}
\usepackage{geometry}
\geometry{margin=2.5cm}

\title{Algebre Relaționale: Reguli, Operatorii și Echivalențe}
\author{Petcu Ionuț Octavian}
\date{}

\begin{document}
\maketitle

\section{Introducere}
Algebra relațională este un limbaj formal pentru manipularea și interogarea datelor în modelul relațional. Ea operează pe relații (tabele) și produce relații noi ca rezultat.

O relație $R$ este un set de tuple:


\[
R \subseteq D_1 \times D_2 \times \dots \times D_n.
\]



\section{Operatorii fundamentali}

\subsection{1. Selecția ($\sigma$)}
\textbf{Definiție:}


\[
\sigma_{\theta}(R) = \{ t \in R \mid t \text{ satisface condiția } \theta \}.
\]



\textbf{Exemplu:}


\[
\sigma_{\text{salariu} > 3000}(Angajati)
\]



\subsection{2. Proiecția ($\pi$)}
\textbf{Definiție:}


\[
\pi_{A_1, A_2, \dots, A_k}(R) = \{ (t[A_1], \dots, t[A_k]) \mid t \in R \}.
\]



\textbf{Exemplu:}


\[
\pi_{nume, salariu}(Angajati)
\]



\subsection{3. Reuniunea ($\cup$)}
\textbf{Definiție:}


\[
R \cup S = \{ t \mid t \in R \text{ sau } t \in S \}.
\]



\textbf{Condiție:} $R$ și $S$ trebuie să fie \textit{compatibile} (aceleași atribute).

\subsection{4. Intersecția ($\cap$)}


\[
R \cap S = \{ t \mid t \in R \text{ și } t \in S \}.
\]



\subsection{5. Diferența ($-$)}


\[
R - S = \{ t \mid t \in R \text{ și } t \notin S \}.
\]



\subsection{6. Produsul cartezian ($\times$)}


\[
R \times S = \{ t_R \circ t_S \mid t_R \in R, t_S \in S \}.
\]



\subsection{7. Renumirea ($\rho$)}


\[
\rho_{A \rightarrow B}(R)
\]



\section{Operatori derivați}

\subsection{1. Join natural ($\bowtie$)}
\textbf{Definiție:}


\[
R \bowtie S = \sigma_{R.A = S.A}(R \times S)
\]


unde $A$ sunt atributele comune.

\textbf{Exemplu:}


\[
Angajati \bowtie Departamente
\]



\subsection{2. Theta-join ($\bowtie_{\theta}$)}


\[
R \bowtie_{\theta} S = \sigma_{\theta}(R \times S)
\]



\subsection{3. Semi-join ($\ltimes$)}


\[
R \ltimes S = \pi_{\text{atribute}(R)}(R \bowtie S)
\]



\subsection{4. Anti-join}


\[
R \setminus\!\!\bowtie S = R - (R \ltimes S)
\]



\subsection{5. Diviziunea ($\div$)}
Pentru relații:


\[
R(X,Y), \quad S(Y)
\]





\[
R \div S = \{ x \mid \forall y \in S, (x,y) \in R \}
\]



\textbf{Echivalent:}


\[
R \div S = \pi_X(R) - \pi_X((\pi_X(R) \times S) - R)
\]



\section{Echivalențe importante}

\subsection{1. Selecția se distribuie peste reuniune}


\[
\sigma_{\theta}(R \cup S) = \sigma_{\theta}(R) \cup \sigma_{\theta}(S)
\]



\subsection{2. Selecția se distribuie peste produs}


\[
\sigma_{\theta}(R \times S) = R \bowtie_{\theta} S
\]



\subsection{3. Proiecția peste selecție}


\[
\pi_L(\sigma_{\theta}(R)) = \sigma_{\theta'}(\pi_L(R))
\]


dacă $\theta$ implică doar atribute din $L$.

\subsection{4. Join natural ca operator derivat}


\[
R \bowtie S = \pi_{A \cup B}(\sigma_{R.C = S.C}(R \times S))
\]



\section{Exemple complete}

\subsection{Exemplul 1: Angajați și departamente}


\[
Rez = \pi_{nume, denumire}(Angajati \bowtie Departamente)
\]



\subsection{Exemplul 2: Angajați cu salariu mare în IT}


\[
Rez = \sigma_{salariu > 5000}(\sigma_{dept = \text{IT}}(Angajati))
\]



\subsection{Exemplul 3: Diviziune}
Găsește angajații care știu \textbf{toate limbajele} din setul $S$.



\[
Rez = Cunostinte(angajat, limbaj) \div S(limbaj)
\]


\section{Reguli fundamentale ale algebrei relaționale}

În această secțiune prezentăm cele mai importante reguli de transformare,
echivalență și optimizare din algebra relațională. Aceste reguli sunt
esențiale pentru rescrierea interogărilor și demonstrarea echivalenței
între expresii relaționale.

\subsection*{Reguli principale}

\begin{enumerate}
    \item \textbf{Selecțiile sunt comutative}
    

\[
        \sigma_{\theta_1}(\sigma_{\theta_2}(R)) = 
        \sigma_{\theta_2}(\sigma_{\theta_1}(R))
    \]



    \item \textbf{Proiecțiile sunt idempotente}
    

\[
        \pi_A(\pi_B(R)) = \pi_A(R) \quad \text{dacă } A \subseteq B
    \]



    \item \textbf{Selecția se distribuie peste reuniune}
    

\[
        \sigma_{\theta}(R \cup S) = 
        \sigma_{\theta}(R) \cup \sigma_{\theta}(S)
    \]



    \item \textbf{Proiecția se distribuie peste reuniune}
    

\[
        \pi_A(R \cup S) = \pi_A(R) \cup \pi_A(S)
    \]



    \item \textbf{Join-ul natural este comutativ}
    

\[
        R \bowtie S = S \bowtie R
    \]



    \item \textbf{Join-ul natural este asociativ}
    

\[
        (R \bowtie S) \bowtie T = R \bowtie (S \bowtie T)
    \]



    \item \textbf{Selecția peste produs cartezian devine theta-join}
    

\[
        \sigma_{\theta}(R \times S) = R \bowtie_{\theta} S
    \]



    \item \textbf{Intersecția exprimată prin diferență}
    

\[
        R \cap S = R - (R - S)
    \]



    \item \textbf{Reuniunea exprimată prin diferență}
    

\[
        R \cup S = (R - S) \cup S
    \]



    \item \textbf{Diviziunea exprimată prin operatori de bază}
    

\[
        R \div S = \pi_X(R) - 
        \pi_X\big((\pi_X(R) \times S) - R\big)
    \]


\end{enumerate}



\end{document}
